\documentclass[12pt,a4paper]{article}
\usepackage{fontspec}
\setmainfont{DejaVu Serif}
\setsansfont{DejaVu Sans}
\setmonofont{DejaVu Sans Mono}[Scale=0.85]
\usepackage[margin=2.2cm]{geometry}
\usepackage{xcolor}
\usepackage{enumitem}
\usepackage{booktabs}
\usepackage{array}
\usepackage{listings}
\usepackage{hyperref}
\usepackage{titlesec}

\titleformat{\section}{\large\bfseries}{\thesection.}{0.5em}{}
\setlength{\parskip}{0.5em}
\setlength{\parindent}{0pt}
\sloppy
\setlength{\emergencystretch}{3em}
\hypersetup{hidelinks}

\lstset{
  basicstyle=\ttfamily\small,
  frame=single, framerule=0.3pt, rulecolor=\color{gray!60},
  backgroundcolor=\color{gray!6},
  columns=fullflexible, keepspaces=true,
  language={}, showstringspaces=false,
  xleftmargin=1em, xrightmargin=1em
}

\begin{document}

\begin{center}
  {\Large\bfseries Лабораторийн ажил №1}\\[0.4em]
  {\large Үйлдлийн систем дээр суурилсан программчлалын орчин бүрдүүлэх}\\[0.8em]
  F.CSM302 Үйлдлийн систем\\
  Гүйцэтгэсэн: Б.Билгүүн (B190910014)\\[0.3em]
  Сурах бичиг: A. Silberschatz, P. B. Galvin, G. Gagne, \textit{Operating System Concepts}, 10th ed., Wiley, 2018 (§1.7, §18)
\end{center}
\hrule
\bigskip

\section{Даалгавар}
\begin{enumerate}[leftmargin=1.5em]
  \item Windows дээр Linux-ийг виртуалаар суулгах, сүлжээ тохируулж виртуал систем рүү файл хуулах.
  \item Виртуал систем рүү SSH хандалт тохируулах (OpenSSH + PuTTY), үндсэн системээс нэвтэрч файлд хандах.
  \item VS Code суулгах.
  \item Үндсэн ба виртуал систем дээр C болон Java хэлээр “Hello world” ажиллуулах.
\end{enumerate}

\section{Бүрдүүлсэн орчин}
\begin{center}
\begin{tabular}{l>{\raggedright}p{5.3cm}>{\raggedright\arraybackslash}p{5.8cm}}
\toprule
 & Үндсэн систем (host) & Виртуал систем (guest)\\
\midrule
ҮС & Windows 11 Home & Ubuntu 22.04.5 LTS (WSL2, кернел 6.6.87.2)\\
C хөрвүүлэгч & gcc 16.2 (WinLibs MinGW-w64) & gcc 11.4 (\texttt{build-essential})\\
Java & JDK 25 & OpenJDK 17\\
Сүлжээ & \texttt{vEthernet (WSL)} 172.24.0.1 & \texttt{eth0} 172.24.x.x (NAT)\\
SSH & PuTTY, OpenSSH клиент & OpenSSH сервер 8.9 (порт 22)\\
Засварлагч & VS Code & —\\
\bottomrule
\end{tabular}
\end{center}

\section{Виртуалчлал — яагаад WSL2 вэ}
Виртуалчлал нь нэг физик компьютер дээр хэд хэдэн үйлдлийн системийг зэрэг ажиллуулах технологи; \textbf{hypervisor} (VMM) техник хангамжийг хийсвэрлэж guest ҮС бүрт “өөрийн компьютер” мэт харагдуулна.
\begin{itemize}[leftmargin=1.5em]
  \item \textbf{Type-1} hypervisor техник хангамж дээр шууд ажиллана (Hyper-V, ESXi); \textbf{Type-2} нь host ҮС дээр программ шиг ажиллана (VirtualBox).
  \item Windows 11 \emph{Home}-д бүрэн Hyper-V байхгүй. \textbf{WSL2} нь хөнгөн VM дотор жинхэнэ Linux кернел ажиллуулдаг тул систем дуудалт, \texttt{fork}, \texttt{strace} бүгд ажиллана; RAM-ыг динамикаар авдаг, Windows-ийн диск \texttt{/mnt/c} дор шууд харагдана.
  \item Ubuntu 22.04 LTS — урт хугацааны дэмжлэгтэй, хамгийн өргөн хэрэглэгддэг дистрибьюшн.
\end{itemize}
\begin{lstlisting}
PS> wsl --install -d Ubuntu-22.04
PS> wsl -l -v                  # Ubuntu-22.04  Running  2   (VERSION 2 = WSL2)
$ uname -a                     # Linux ... microsoft-standard-WSL2 x86_64
\end{lstlisting}

\section{Сүлжээ ба файл хуулах}
WSL2 нь \textbf{NAT} горимтой: Windows дээр \texttt{vEthernet (WSL)} адаптер (172.24.0.1) үүсч, VM 172.24.x.x хаяг авна; гадагш host-ийн хаягаар гарна.
\begin{lstlisting}
$ ip addr show eth0 | grep inet      # VM-ийн IP
$ ping -c 2 172.24.0.1               # VM -> Windows
$ ping -c 2 8.8.8.8                  # VM -> интернэт
$ cp "/mnt/c/lesson 2026/F.CSA302/lab1/test.txt" ~/ && cat ~/test.txt
Hello from Windows
\end{lstlisting}
Файл хуулах өөр хоёр арга: File Explorer-т \texttt{\textbackslash\textbackslash wsl\$\textbackslash Ubuntu-22.04} хаягаар чирэх, SSH-ээр \texttt{pscp}/\texttt{scp}.

\section{SSH хандалт}
SSH нь клиент–сервер, шифрлэгдсэн протокол (порт 22). Нууц үгийн оронд \textbf{ed25519 түлхүүр хос} ашигласан: public key VM-ийн \texttt{\textasciitilde/.ssh/authorized\_keys}-д, private key зөвхөн Windows дээр (PuTTY-д \texttt{.ppk} формат). Сервер клиентээс private key-ээр гарын үсэг зурахыг шаардаж шалгадаг тул нууц үг сүлжээгээр дамжихгүй.
\begin{lstlisting}
$ apt install openssh-server && systemctl enable --now ssh
$ ss -tlnp | grep :22                       # 0.0.0.0:22 дээр сонсож байна
PS> ssh-keygen -t ed25519 -f ~/.ssh/wsl_lab1 # түлхүүр хос
PuTTY: Host = localhost, Port = 22, user = root, key = wsl_lab1.ppk
$ whoami; ls "/mnt/c/lesson 2026/F.CSA302"   # нэвтэрч файлд хандав
\end{lstlisting}

\section{Hello world — хоёр систем дээр}
\begin{lstlisting}
PS> gcc hello.c -o hello.exe; .\hello.exe        # Windows: PE (.exe)
Hello world
PS> javac HelloWorld.java; java HelloWorld
Hello world
$ gcc hello.c -o hello && ./hello                # Ubuntu: ELF
Hello world
$ java HelloWorld.java
Hello world
$ file hello        # ELF 64-bit LSB executable, x86-64
\end{lstlisting}
C код \texttt{gcc}-ээр тухайн ҮС/CPU-ийн \textbf{машины код} болдог (Windows — PE, Linux — ELF) тул ҮС бүр дээр тусад нь хөрвүүлнэ. Java код \texttt{javac}-аар \textbf{bytecode} (\texttt{.class}) болж JVM дээр ажилладаг тул нэг \texttt{.class} хоёр ҮС дээр адилхан ажиллана. \texttt{printf} нь Linux дээр эцэстээ \texttt{write(1, ...)} систем дуудалт хийдэг (\texttt{strace ./hello}).

\section{Дүгнэлт}
Windows 11 дээр WSL2-оор Ubuntu 22.04 виртуал систем суулгаж, NAT сүлжээ, файл хуулалт, түлхүүрээр нэвтрэх SSH, VS Code-ийг тохируулж, C ба Java программыг хоёр систем дээр хөрвүүлэн ажиллуулав. Энэ орчныг дараагийн бүх лабд (gcc, strace, pthreads, XeLaTeX) ашигласан.

\end{document}
